% Options for packages loaded elsewhere
\PassOptionsToPackage{unicode}{hyperref}
\PassOptionsToPackage{hyphens}{url}
\documentclass[
]{article}
\usepackage{xcolor}
\usepackage{amsmath,amssymb}
\setcounter{secnumdepth}{-\maxdimen} % remove section numbering
\usepackage{iftex}
\ifPDFTeX
  \usepackage[T1]{fontenc}
  \usepackage[utf8]{inputenc}
  \usepackage{textcomp} % provide euro and other symbols
\else % if luatex or xetex
  \usepackage{unicode-math} % this also loads fontspec
  \defaultfontfeatures{Scale=MatchLowercase}
  \defaultfontfeatures[\rmfamily]{Ligatures=TeX,Scale=1}
\fi
\usepackage{lmodern}
\ifPDFTeX\else
  % xetex/luatex font selection
\fi
% Use upquote if available, for straight quotes in verbatim environments
\IfFileExists{upquote.sty}{\usepackage{upquote}}{}
\IfFileExists{microtype.sty}{% use microtype if available
  \usepackage[]{microtype}
  \UseMicrotypeSet[protrusion]{basicmath} % disable protrusion for tt fonts
}{}
\makeatletter
\@ifundefined{KOMAClassName}{% if non-KOMA class
  \IfFileExists{parskip.sty}{%
    \usepackage{parskip}
  }{% else
    \setlength{\parindent}{0pt}
    \setlength{\parskip}{6pt plus 2pt minus 1pt}}
}{% if KOMA class
  \KOMAoptions{parskip=half}}
\makeatother
\usepackage{color}
\usepackage{fancyvrb}
\newcommand{\VerbBar}{|}
\newcommand{\VERB}{\Verb[commandchars=\\\{\}]}
\DefineVerbatimEnvironment{Highlighting}{Verbatim}{commandchars=\\\{\}}
% Add ',fontsize=\small' for more characters per line
\newenvironment{Shaded}{}{}
\newcommand{\AlertTok}[1]{\textcolor[rgb]{1.00,0.00,0.00}{\textbf{#1}}}
\newcommand{\AnnotationTok}[1]{\textcolor[rgb]{0.38,0.63,0.69}{\textbf{\textit{#1}}}}
\newcommand{\AttributeTok}[1]{\textcolor[rgb]{0.49,0.56,0.16}{#1}}
\newcommand{\BaseNTok}[1]{\textcolor[rgb]{0.25,0.63,0.44}{#1}}
\newcommand{\BuiltInTok}[1]{\textcolor[rgb]{0.00,0.50,0.00}{#1}}
\newcommand{\CharTok}[1]{\textcolor[rgb]{0.25,0.44,0.63}{#1}}
\newcommand{\CommentTok}[1]{\textcolor[rgb]{0.38,0.63,0.69}{\textit{#1}}}
\newcommand{\CommentVarTok}[1]{\textcolor[rgb]{0.38,0.63,0.69}{\textbf{\textit{#1}}}}
\newcommand{\ConstantTok}[1]{\textcolor[rgb]{0.53,0.00,0.00}{#1}}
\newcommand{\ControlFlowTok}[1]{\textcolor[rgb]{0.00,0.44,0.13}{\textbf{#1}}}
\newcommand{\DataTypeTok}[1]{\textcolor[rgb]{0.56,0.13,0.00}{#1}}
\newcommand{\DecValTok}[1]{\textcolor[rgb]{0.25,0.63,0.44}{#1}}
\newcommand{\DocumentationTok}[1]{\textcolor[rgb]{0.73,0.13,0.13}{\textit{#1}}}
\newcommand{\ErrorTok}[1]{\textcolor[rgb]{1.00,0.00,0.00}{\textbf{#1}}}
\newcommand{\ExtensionTok}[1]{#1}
\newcommand{\FloatTok}[1]{\textcolor[rgb]{0.25,0.63,0.44}{#1}}
\newcommand{\FunctionTok}[1]{\textcolor[rgb]{0.02,0.16,0.49}{#1}}
\newcommand{\ImportTok}[1]{\textcolor[rgb]{0.00,0.50,0.00}{\textbf{#1}}}
\newcommand{\InformationTok}[1]{\textcolor[rgb]{0.38,0.63,0.69}{\textbf{\textit{#1}}}}
\newcommand{\KeywordTok}[1]{\textcolor[rgb]{0.00,0.44,0.13}{\textbf{#1}}}
\newcommand{\NormalTok}[1]{#1}
\newcommand{\OperatorTok}[1]{\textcolor[rgb]{0.40,0.40,0.40}{#1}}
\newcommand{\OtherTok}[1]{\textcolor[rgb]{0.00,0.44,0.13}{#1}}
\newcommand{\PreprocessorTok}[1]{\textcolor[rgb]{0.74,0.48,0.00}{#1}}
\newcommand{\RegionMarkerTok}[1]{#1}
\newcommand{\SpecialCharTok}[1]{\textcolor[rgb]{0.25,0.44,0.63}{#1}}
\newcommand{\SpecialStringTok}[1]{\textcolor[rgb]{0.73,0.40,0.53}{#1}}
\newcommand{\StringTok}[1]{\textcolor[rgb]{0.25,0.44,0.63}{#1}}
\newcommand{\VariableTok}[1]{\textcolor[rgb]{0.10,0.09,0.49}{#1}}
\newcommand{\VerbatimStringTok}[1]{\textcolor[rgb]{0.25,0.44,0.63}{#1}}
\newcommand{\WarningTok}[1]{\textcolor[rgb]{0.38,0.63,0.69}{\textbf{\textit{#1}}}}
\usepackage{longtable,booktabs,array}
\usepackage{calc} % for calculating minipage widths
% Correct order of tables after \paragraph or \subparagraph
\usepackage{etoolbox}
\makeatletter
\patchcmd\longtable{\par}{\if@noskipsec\mbox{}\fi\par}{}{}
\makeatother
% Allow footnotes in longtable head/foot
\IfFileExists{footnotehyper.sty}{\usepackage{footnotehyper}}{\usepackage{footnote}}
\makesavenoteenv{longtable}
\setlength{\emergencystretch}{3em} % prevent overfull lines
\providecommand{\tightlist}{%
  \setlength{\itemsep}{0pt}\setlength{\parskip}{0pt}}
\usepackage{bookmark}
\IfFileExists{xurl.sty}{\usepackage{xurl}}{} % add URL line breaks if available
\urlstyle{same}
\hypersetup{
  hidelinks,
  pdfcreator={LaTeX via pandoc}}

\author{}
\date{}

\begin{document}

\section{A Power-Law Cost Model for Operational-Risk Incident Reports:
Calibration on 40 Historical Incidents and Implications for DORA
Compliance}\label{a-power-law-cost-model-for-operational-risk-incident-reports-calibration-on-40-historical-incidents-and-implications-for-dora-compliance}

\textbf{Alberto Giovanni Gerli}\^{}\{1,2\}

\^{}1 Tourbillon Tech Srl, Padova, Italy \^{}2 Dipartimento di Scienze
Cliniche e di Comunità, Università degli Studi di Milano, Italy

\textbf{Corresponding author:} alberto@albertogerli.it

\textbf{Keywords:} operational risk, DORA, EU Regulation 2022/2554,
major-incident reporting, power-law cost model, bootstrap inference,
regime-switching HMM, leave-one-out cross-validation

\begin{center}\rule{0.5\linewidth}{0.5pt}\end{center}

\subsection{Abstract}\label{abstract}

The Digital Operational Resilience Act (Regulation EU 2022/2554,
``DORA'') requires regulated financial entities to report the economic
impact of major information-and-communication-technology incidents
within deadlines defined by a seven-criterion classification. Reporting
templates request a euro figure or band for
\texttt{economic\_impact\_eur} even when actuals are not yet known and
the reporting institution must rely on a prospective estimate. Existing
approaches anchor the estimate to either historical-cost analogues
selected by hand or to qualitative tier mappings prescribed by internal
risk-management policy. Both approaches lack a defensible quantitative
model and a transparent uncertainty band that survives external audit.

We propose a calibrated power-law cost model
\(\hat{c}(s) = \beta \cdot s^\gamma\) where \(s\) is a scalar
shock-magnitude index of the incident and \((\beta, \gamma)\) are
estimated by log-log Huber regression on a curated reference table of
\(N=40\) historical incidents (1998--2024) across seven operational
categories --- banking-IT, banking-EU, banking-US, sovereign, cyber,
telco, and energy. We document the dataset, the per-category fits, and
the model selection rule that promotes the power-law from diagnostic to
primary estimator only when the log-space coefficient of determination
exceeds \(0.5\).

A leave-one-out cross-validation reports a hit-rate within \(\pm 100\%\)
of \(80\%\) under the power-law model, against \(35\%\) for a linear
baseline \(\hat{c}(s) = \alpha s\) and \(40\%\) for a per-category
linear baseline. Median absolute percent error drops from \(394\%\) to
\(57\%\). The remaining error concentrates in the banking-US bucket
where Lehman Brothers (2008) is a high-leverage outlier dominating the
in-bucket slope.

The headline estimate is reported alongside six independent statistical
diagnostics: (i) an empirical pairs-bootstrap on \((\beta, \gamma)\) for
the epistemic uncertainty interval, (ii) the
Eicker--Huber--White--MacKinnon HC3 sandwich standard error, (iii) a
2-state Gaussian hidden Markov model on monthly log(VIX) for
incident-level regime posteriors and a regime-mixture α as
cross-validation, (iv) a Hill estimator for the Pareto tail index of
residuals to flag infinite-variance regimes, (v) two-stage least squares
with the regime posterior as instrument for the shock index to test for
endogeneity, and (vi) the log-log slope reported as a Taleb-style
fragility exponent. The diagnostic stack is designed so that the
supervisor reads each output against a falsifiable assumption, rather
than against a black-box estimator.

We release the reference dataset, the calibration code, and the
validation harness as an open benchmark for operational-risk cost
modeling. The benchmark provides a reproducible skill floor for proposed
cost estimators and a curated dataset for empirical work in
operational-risk modeling under DORA and adjacent regulatory frameworks.

\begin{center}\rule{0.5\linewidth}{0.5pt}\end{center}

\subsection{1. Introduction}\label{introduction}

\subsubsection{1.1 Motivation}\label{motivation}

Article 19 of Regulation EU 2022/2554 (DORA, in force since January
2025) requires regulated financial entities to submit major-incident
reports to their competent supervisory authority --- for Italian banks
the Banca d'Italia, for insurance undertakings IVASS, and for harmonised
cross-jurisdictional cases the European Banking Authority, EIOPA, or
ESMA depending on the entity. Annex I of the implementing technical
standards (RTS, Joint Committee Final Report JC 2024-43) defines a
seven-criterion classification matrix that includes a quantitative
\texttt{economic\_impact\_eur} field. The matrix prescribes four
reporting tiers --- non-material (\textless€100k), low (€100k--€1m),
high (€1m--€10m), critical (\(\geq\) €10m) --- and the entity is
required to select a tier within four hours of an event being classified
as major.

In practice the reporting institution must form a prospective estimate
of the economic impact within a short deadline and based on partial
information, and revise that estimate as the incident evolves. Neither
the regulation nor the technical standards prescribe a specific
estimation methodology, leaving the choice to the institution's internal
risk-management policy. Existing implementations fall into two families.

The first family is \textbf{historical-cost analogue selection}. The
risk officer identifies one or two past incidents that resemble the
current one --- typically by the same institution or by a comparable
peer --- and reports the analogue's eventual cost. The approach is
intuitive and survives audit when the analogue is well-documented, but
suffers from selection bias and provides no formal uncertainty band: a
single analogue is a sample of size one. Industry guidance from the
Operational Risk Management Forum (ORMF, 2021) and the EBA Guidelines on
ICT and security risk management (EBA/GL/2019/04) acknowledges the
practice but does not endorse it for the DORA reporting context.

The second family is \textbf{qualitative tier mapping}. The institution
defines internal heuristics that map incident attributes (number of
clients affected, downtime hours, geographical spread) to one of the
four DORA tiers, typically through a rule-based engine implemented in
the institution's GRC platform (MetricStream, ServiceNow GRC, IBM
OpenPages, Wolters Kluwer OneSumX). The mapping is auditable but the
tier boundaries are coarse: an incident estimated at €5m and one
estimated at €15m fall into different tiers and trigger different
supervisory follow-ups, but the underlying analysis cannot distinguish
between them.

Neither approach provides what an external auditor or the supervisory
authority would expect of a quantitative model: a documented
relationship between observable inputs and the cost estimate, a
calibrated parameter set with measurable goodness-of-fit on historical
data, and a propagated uncertainty band with formal coverage properties.

\subsubsection{1.2 Contribution}\label{contribution}

We close this gap by proposing a \textbf{calibrated power-law cost
model} for operational-risk economic-impact estimation:

\[
\hat{c}(s) = \beta \cdot s^\gamma + \varepsilon, \qquad \log \hat{c} = \log \beta + \gamma \log s + u,
\]

where \(s\) is a scalar shock-magnitude index of the incident,
\((\beta, \gamma)\) are category-specific parameters fitted by log-log
Huber regression on a curated reference table of \(N=40\) historical
incidents (1998--2024), and \(\varepsilon\) has empirical distribution
estimated from a pairs-bootstrap on the calibration subset.

The contributions are:

\begin{enumerate}
\def\labelenumi{\arabic{enumi}.}
\item
  \textbf{A curated open reference dataset}
  (\texttt{shared/dora\_reference\_incidents.json}, schema v1.1) of
  \(N=40\) historical incidents across seven operational categories,
  with per-incident shock-magnitude estimates, public-domain euro costs,
  ISO incident dates, regime labels, and a minimum of two source
  citations per row. The dataset is auditable, version-controlled, and
  released as the basis for an open benchmark for operational-risk cost
  modeling (Section 3).
\item
  \textbf{Empirical evidence that the cost-vs-shock relationship is
  super-linear with \(\hat{\gamma} \approx 3\) on the overall corpus}
  (\(R^2_{\log} = 0.72\)) and \(\hat{\gamma} = 3.12\) on the banking-IT
  slice (\(R^2_{\log} = 0.97\)). The linear cost model
  \(\hat{c}(s) = \alpha s\), which is implicit in industry practice, is
  therefore mis-specified: it systematically over-predicts small
  incidents and under-predicts large ones (Section 4).
\item
  \textbf{A model selection rule} that promotes the power-law from
  diagnostic to primary estimator when the per-category log-space
  \(R^2 \geq 0.5\), and falls back to the linear baseline otherwise. The
  rule is data-driven, audit-friendly (it produces an interpretable
  label per request), and conservative (smaller and noisier buckets
  default to the more familiar baseline) (Section 4.4).
\item
  \textbf{A leave-one-out cross-validation under three competing model
  specifications} (overall linear, per-category linear, per-category
  power-law) on the reference dataset. The power-law specification
  attains \(80\%\) within \(\pm 100\%\) against \(40\%\) for the
  per-category linear baseline and \(35\%\) for the overall linear
  baseline; median absolute percent error drops from \(394\%\) to
  \(57\%\) (Section 6).
\item
  \textbf{A six-layer diagnostic stack} alongside the headline estimate:
  empirical pairs-bootstrap on \((\beta, \gamma)\), HC3
  heteroscedastic-robust standard error, hidden Markov regime
  posteriors, Hill tail-index estimator, two-stage least squares with
  regime as instrument, and a log-log slope reported as a fragility
  exponent. Each diagnostic targets a specific assumption underlying the
  headline estimate, and each is designed to be falsifiable and
  externally auditable (Section 5).
\item
  \textbf{An open benchmark and reproducibility package} that allows any
  proposed cost estimator to be evaluated against the same dataset under
  the same leave-one-out protocol, providing a fixed reference skill
  floor for the operational-risk modeling literature (Section 8).
\end{enumerate}

\subsubsection{1.3 Paper organisation}\label{paper-organisation}

Section 2 reviews related work in operational-risk modeling, power-law
regularities in financial losses, regime-switching econometrics, and the
regulatory landscape under DORA. Section 3 documents the reference
incident dataset, its construction methodology, and aggregate
statistics. Section 4 presents the linear and power-law specifications
and the model selection rule. Section 5 describes the six-layer
diagnostic stack. Section 6 presents the cross-validation results.
Section 7 discusses limitations and external validity. Section 8
describes the open-benchmark release. Section 9 concludes.

\begin{center}\rule{0.5\linewidth}{0.5pt}\end{center}

\subsection{2. Related Work}\label{related-work}

\subsubsection{2.1 Operational risk
modeling}\label{operational-risk-modeling}

The quantitative literature on operational risk historically separated
\emph{frequency} models (how often does a loss occur) from
\emph{severity} models (how large is each loss). The Loss Distribution
Approach (LDA), formalised under the Basel II Advanced Measurement
Approach in 2004 and refined under Basel III, fits a Poisson process to
event frequency and a heavy-tailed parametric distribution (typically
lognormal, Weibull, or generalised Pareto) to event severity, then
convolves the two via Monte Carlo to estimate the 99.9\% Value-at-Risk
one-year aggregate loss {[}Frachot et al.~2001; Cruz 2002; Aue \&
Kalkbrener 2006; Chernobai et al.~2007{]}.

The LDA framework is appropriate for \emph{aggregate-portfolio} risk
capital but is less well-suited to \emph{single-incident} prospective
cost estimation, which is what DORA Article 19 requires. The institution
observes a specific event in progress, is asked to estimate its cost
within hours, and must condition on the incident's category, geography,
and counterparty exposure rather than on a portfolio-level
frequency--severity calibration.

Recent work has begun to address this gap. Cope et al.~(2009) document
that operational losses across institutions exhibit strong scaling
regularities by event type, suggesting that conditional models
stratified by category outperform pooled fits. Hess (2011) analyses the
Algo FIRST loss database and reports that within-category severity
distributions are well-approximated by power-law tails. Bardoscia et
al.~(2017) propose a network-conditional framework for operational
losses in banking. None of these target the specific DORA major-incident
reporting use-case, and none provide a public benchmark.

\subsubsection{2.2 Power laws in financial
losses}\label{power-laws-in-financial-losses}

The empirical literature on financial losses has consistently reported
power-law (Pareto-type) tails. Gabaix et al.~(2003) documented power-law
tails in stock-market returns; Bouchaud et al.~(2002) found similar
regularities in volatility and trading volume. For
systemic-banking-crisis losses specifically, Reinhart and Rogoff (2009)
and Laeven and Valencia (2018) maintain a curated dataset whose cost
distribution exhibits a heavy right tail consistent with power-law
behaviour above the median.

Power-law cost relationships of the form \(c \propto s^\gamma\) with
\(\gamma > 1\) have been theoretically derived in cascade and contagion
models (Acemoglu et al.~2015; Battiston et al.~2012) and empirically
confirmed in network-banking simulations (Caccioli et al.~2014). The
\(\gamma\) exponent is variously called \emph{contagion intensity},
\emph{fragility exponent} {[}Taleb 2012{]}, or simply \emph{cost
convexity}. Our work applies this regularity to the DORA reporting
context and validates it against a held-out historical sample for the
first time, to our knowledge.

\subsubsection{2.3 Regime-switching
econometrics}\label{regime-switching-econometrics}

The 2-state Gaussian hidden Markov model for financial-market regime
identification is due to Hamilton (1989, \emph{Econometrica} 57:357),
with EM estimation following Baum--Welch (Baum et al.~1970) and the
standard log-space implementation (Rabiner 1989). Applications to
volatility regime detection in equity markets (Schwert 1989; Pagan \&
Schwert 1990; Ang \& Bekaert 2002), credit-spread regime classification
(Kiefer 2010), and central-bank policy phase identification (Sims \& Zha
2006) are extensive.

We use Hamilton's 2-state regime-switching HMM as a \emph{diagnostic}
layer: the regime posterior conditioned on incident date is one of three
independent stress signals (the other two being category and hand-coded
regime label), and the regime-mixture α serves as a cross-check on the
per-category fit. This is a non-standard application --- the literature
uses HMM regime posteriors as inputs to a forecasting or pricing model,
not as cross-validation for a separately fitted cost surface. The use is
closer in spirit to Ang \& Timmermann (2012, \emph{Annual Review of
Financial Economics} 4:313) who treat regime classification as an
exploratory device for empirical-finance datasets.

\subsubsection{2.4 Robust regression and bootstrap
inference}\label{robust-regression-and-bootstrap-inference}

The Huber M-estimator (Huber 1964) is the standard robust regression
specification for samples containing high-leverage outliers; we use it
in both the level (linear \(\alpha s\)) and log-space (power-law) fits
with the canonical tuning constant \(k = 1.345\). The empirical
pairs-bootstrap follows Efron (1979) and Davison \& Hinkley (1997,
\emph{Bootstrap Methods and Their Application}); we report empirical
5°/95° quantiles of the resampled coefficient distribution rather than
parametric confidence intervals because the residual distribution is
heavy-tailed and a Gaussian band would understate tail risk.

The Eicker--Huber--White--MacKinnon HC3 sandwich estimator (Eicker 1967;
Huber 1967; White 1980; MacKinnon \& White 1985) provides
heteroscedastic-robust standard errors without requiring the analyst to
specify a variance form. The Hill estimator for the Pareto tail index
(Hill 1975, \emph{Annals of Statistics} 3:1163) and its small-sample
bias properties (Hall 1990; Resnick \& Stărică 1997) are standard tools
in extreme-value analysis; we use it as a diagnostic on residual tails
to flag the infinite-variance regime that occurs when the corpus
contains a Lehman-class outlier.

\subsubsection{2.5 The DORA regulatory
context}\label{the-dora-regulatory-context}

Regulation EU 2022/2554 (DORA) entered into force on 16 January 2023 and
is applicable from 17 January 2025. Article 17 requires regulated
entities to maintain an ICT-related incident management process; Article
18 establishes the criteria for classifying ICT-related incidents as
``major''; Article 19 prescribes the reporting obligation to the
competent supervisory authority. The implementing technical standards
were finalised in the EBA/EIOPA/ESMA Joint Committee Final Report JC
2024-43 (July 2024) and define the seven-criterion classification matrix
(clients\_affected, data\_losses, reputational\_impact,
duration\_downtime\_hours, geographical\_spread, economic\_impact\_eur,
criticality\_of\_services\_affected) along with the report templates.

Existing surveys of the DORA implementation landscape (Deloitte 2024;
PwC 2024; Accenture 2025) emphasise process and governance compliance
and largely defer the quantitative \texttt{economic\_impact\_eur}
estimation to existing institutional GRC tooling. Our work targets that
specific quantitative gap.

\begin{center}\rule{0.5\linewidth}{0.5pt}\end{center}

\subsection{3. Reference Incident
Dataset}\label{reference-incident-dataset}

\subsubsection{3.1 Construction
methodology}\label{construction-methodology}

The reference dataset, \texttt{shared/dora\_reference\_incidents.json}
(schema v1.1), comprises \(N=40\) historical operational-risk incidents
in regulated financial-system or financial-system-adjacent sectors that
occurred between September 1998 and July 2024. Each incident is an event
for which (i) a public-domain estimate of the eventual euro cost exists
in regulator filings, financial-stability reports, peer-reviewed
retrospectives, or official institutional disclosures, (ii) a rough
estimate of the incident's shock magnitude can be derived from
contemporaneous news intensity, and (iii) the event date and category
are unambiguous.

The construction proceeded in three phases. First, an initial seed list
of \(\sim 80\) candidate events was assembled from the Reinhart-Rogoff
systemic banking crisis database, Laeven-Valencia 2018, the IMF
Financial Soundness Indicators incident registry, the BIS Quarterly
Review historical archive (2007--2024), and the EBA risk dashboards.
Second, candidates were filtered to those with at least two independent
public-domain cost citations and a primary regulatory or central-bank
source. Third, each surviving candidate was assigned a categorical label
(one of seven operational categories), an ISO incident date for the peak
news event (typically the day of formal announcement: Lehman bankruptcy
filing 2008-09-15, MPS bailout 2017-07-04, CrowdStrike Falcon outage
2024-07-19), and a shock-magnitude estimate following the heuristic
tiers defined in Section 3.3.

The final corpus contains \(N=40\) entries distributed across the seven
categories as follows: banking-IT (Italian banks, \(n=6\)), banking-EU
(\(n=8\)), banking-US (\(n=7\)), sovereign (\(n=6\)), cyber (\(n=6\)),
telco (\(n=4\)), and energy (\(n=3\)). The category distribution
reflects both the documented public-information availability and the
relative importance of each category to the EU regulated-financial
perimeter that DORA covers.

\subsubsection{3.2 Data schema}\label{data-schema}

Each incident row is a JSON object with the following fields:

\begin{Shaded}
\begin{Highlighting}[]
\FunctionTok{\{}
  \DataTypeTok{"id"}\FunctionTok{:} \StringTok{"mps\_bailin\_2017"}\FunctionTok{,}
  \DataTypeTok{"shock\_units"}\FunctionTok{:} \FloatTok{1.6}\FunctionTok{,}
  \DataTypeTok{"cost\_eur\_m"}\FunctionTok{:} \DecValTok{3900}\FunctionTok{,}
  \DataTypeTok{"category"}\FunctionTok{:} \StringTok{"banking\_it"}\FunctionTok{,}
  \DataTypeTok{"label"}\FunctionTok{:} \StringTok{"MPS precautionary recapitalisation (2017)"}\FunctionTok{,}
  \DataTypeTok{"incident\_date"}\FunctionTok{:} \StringTok{"2017{-}07{-}04"}\FunctionTok{,}
  \DataTypeTok{"regime"}\FunctionTok{:} \StringTok{"stressed"}\FunctionTok{,}
  \DataTypeTok{"sources"}\FunctionTok{:} \OtherTok{[}
    \StringTok{"DG Comp state aid SA.47677"}\OtherTok{,}
    \StringTok{"MEF press release 2017{-}07{-}04"}
  \OtherTok{]}
\FunctionTok{\}}
\end{Highlighting}
\end{Shaded}

The fields and their interpretation:

\begin{itemize}
\tightlist
\item
  \texttt{id}: a snake-case identifier, primary key.
\item
  \texttt{shock\_units}: a scalar in the open range \((0, 5)\) measuring
  the magnitude of the event on a heuristic four-tier scale (Section
  3.3).
\item
  \texttt{cost\_eur\_m}: the public-domain estimate of the eventual cost
  in euro-millions, taken from the cited primary source. For events
  denominated in non-euro currencies (Lehman, SVB, Argentina default),
  the conversion uses the historical average exchange rate over the
  resolution window.
\item
  \texttt{category}: one of \texttt{banking\_it}, \texttt{banking\_eu},
  \texttt{banking\_us}, \texttt{sovereign}, \texttt{cyber},
  \texttt{telco}, \texttt{energy}.
\item
  \texttt{label}: a human-readable description used in tables and
  visualisations.
\item
  \texttt{incident\_date}: the ISO date of the peak event (formal
  announcement, takeover, declaration). Used in Section 5 for the
  regime-posterior diagnostic.
\item
  \texttt{regime}: one of \texttt{calm}, \texttt{stressed},
  \texttt{crisis}. Hand-coded from market context. Used as a
  hand-curated cross-check against the data-driven regime posterior in
  Section 5.
\item
  \texttt{sources}: a list of at least two public-domain references.
  Allows audit of the cost figure and fact-checking of the date and
  category.
\end{itemize}

\subsubsection{3.3 Shock-magnitude
heuristic}\label{shock-magnitude-heuristic}

The shock-magnitude index \(s\) is a scalar that aggregates the event's
intrinsic severity and its observed news intensity. The index is
dimensionless and is calibrated to four heuristic tiers reflecting the
scale at which institutional risk officers typically reason about
operational events:

\begin{longtable}[]{@{}
  >{\raggedright\arraybackslash}p{(\linewidth - 6\tabcolsep) * \real{0.2500}}
  >{\raggedright\arraybackslash}p{(\linewidth - 6\tabcolsep) * \real{0.2500}}
  >{\raggedright\arraybackslash}p{(\linewidth - 6\tabcolsep) * \real{0.2500}}
  >{\raggedright\arraybackslash}p{(\linewidth - 6\tabcolsep) * \real{0.2500}}@{}}
\toprule\noalign{}
\begin{minipage}[b]{\linewidth}\raggedright
Tier
\end{minipage} & \begin{minipage}[b]{\linewidth}\raggedright
Range
\end{minipage} & \begin{minipage}[b]{\linewidth}\raggedright
Description
\end{minipage} & \begin{minipage}[b]{\linewidth}\raggedright
Examples
\end{minipage} \\
\midrule\noalign{}
\endhead
\bottomrule\noalign{}
\endlastfoot
I & 0.5 -- 1.2 & Single-firm contained event & Tercas FITD intervention
(2014, \(s=0.8\)); Popolare Bari rescue (2019, \(s=1.0\)) \\
II & 1.2 -- 2.0 & Sector-wide moderate or protracted single-firm & MPS
deposit run (2016, \(s=1.4\)); Wirecard collapse (2020, \(s=1.6\)) \\
III & 2.0 -- 3.0 & Major sector or sovereign-adjacent & SVB collapse
(2023, \(s=2.4\)); Cyprus bail-in (2013, \(s=2.3\)) \\
IV & \(\geq 3.0\) & Systemic or sovereign-class & Brexit Wave-1 (2016,
\(s=3.2\)); Argentina default (2001, \(s=3.5\)); Lehman (2008,
\(s=4.0\)) \\
\end{longtable}

The tier assignment was performed by the author with reference to
contemporaneous press archives, the cited primary regulatory sources,
and the central-bank macro-financial reviews from the relevant year. The
heuristic is intentionally coarse (one significant figure per
assignment) to avoid over-fitting the cost model to subjective
shock-magnitude refinements. Section 7.2 discusses the sensitivity of
the headline cost estimate to plausible shock-magnitude perturbations.

\subsubsection{3.4 Aggregate dataset
statistics}\label{aggregate-dataset-statistics}

Across the \(N=40\) corpus, the unweighted distribution of
shock-magnitude is centred around \(\bar{s} = 1.93\) with standard
deviation \(0.78\) and range \([0.7, 4.0]\). The unweighted distribution
of euro cost is heavily right-skewed: median cost €5,200m, mean cost
€31,860m, maximum cost €600,000m (Lehman). The cost distribution
exhibits a Gini coefficient of \(0.81\), indicating concentration of
mass in a few high-leverage observations. Per-category statistics are
reported in Table 1.

\textbf{Table 1: Reference dataset summary statistics by category.}

\begin{longtable}[]{@{}
  >{\raggedright\arraybackslash}p{(\linewidth - 10\tabcolsep) * \real{0.1667}}
  >{\raggedright\arraybackslash}p{(\linewidth - 10\tabcolsep) * \real{0.1667}}
  >{\raggedright\arraybackslash}p{(\linewidth - 10\tabcolsep) * \real{0.1667}}
  >{\raggedright\arraybackslash}p{(\linewidth - 10\tabcolsep) * \real{0.1667}}
  >{\raggedright\arraybackslash}p{(\linewidth - 10\tabcolsep) * \real{0.1667}}
  >{\raggedright\arraybackslash}p{(\linewidth - 10\tabcolsep) * \real{0.1667}}@{}}
\toprule\noalign{}
\begin{minipage}[b]{\linewidth}\raggedright
Category
\end{minipage} & \begin{minipage}[b]{\linewidth}\raggedright
\(n\)
\end{minipage} & \begin{minipage}[b]{\linewidth}\raggedright
mean \(s\)
\end{minipage} & \begin{minipage}[b]{\linewidth}\raggedright
mean \(c\) (€M)
\end{minipage} & \begin{minipage}[b]{\linewidth}\raggedright
median \(c\) (€M)
\end{minipage} & \begin{minipage}[b]{\linewidth}\raggedright
max \(c\) (€M)
\end{minipage} \\
\midrule\noalign{}
\endhead
\bottomrule\noalign{}
\endlastfoot
banking\_it & 6 & 1.18 & 2,300 & 1,950 & 5,200 \\
banking\_eu & 8 & 1.81 & 21,250 & 7,150 & 90,000 \\
banking\_us & 7 & 2.27 & 96,500 & 9,000 & 600,000 \\
sovereign & 6 & 2.57 & 35,330 & 27,500 & 82,000 \\
cyber & 6 & 2.25 & 21,800 & 7,000 & 100,000 \\
telco & 4 & 1.05 & 3,725 & 1,600 & 11,000 \\
energy & 3 & 1.40 & 15,130 & 11,000 & 34,000 \\
\end{longtable}

The banking-US bucket is dominated by Lehman Brothers (2008, \(c =\)
€600B, \(s = 4.0\)), which is approximately \(66\times\) larger than the
next-largest in-category observation (Washington Mutual, \(c =\) €25B).
Retaining Lehman in the calibration set ensures the model can
extrapolate to systemic-class events, but it dominates per-category fits
in banking-US through high leverage.

The full corpus is reproduced in Appendix A; the JSON file is available
at the project repository (Section 8).

\begin{center}\rule{0.5\linewidth}{0.5pt}\end{center}

\subsection{4. Power-Law Cost Model}\label{power-law-cost-model}

\subsubsection{4.1 Model specifications}\label{model-specifications}

We compare three model specifications fitted to the reference dataset:

\textbf{Model M1 --- Linear pooled.} The simplest specification: a
single slope \(\alpha\) across all incidents, no intercept, no category
conditioning.

\[
M_1: \quad c_i = \alpha \cdot s_i + \varepsilon_i, \quad i = 1, \ldots, N.
\]

\textbf{Model M2 --- Linear per-category.} A category-specific slope
\(\alpha_k\) for each of the seven operational categories, no intercept.

\[
M_2: \quad c_i = \alpha_{k(i)} \cdot s_i + \varepsilon_i, \quad k(i) \in \{1, \ldots, 7\}.
\]

\textbf{Model M3 --- Power-law per-category (proposed).} A
category-specific multiplicative pre-factor \(\beta_k\) and a
category-specific exponent \(\gamma_k\).

\[
M_3: \quad c_i = \beta_{k(i)} \cdot s_i^{\gamma_{k(i)}} + \varepsilon_i,
\qquad \log c_i = \log \beta_{k(i)} + \gamma_{k(i)} \log s_i + u_i.
\]

The power-law model M3 is fitted by ordinary least squares in log space.
Models M1 and M2 are fitted by Huber regression (Huber 1964) with the
canonical tuning constant \(k = 1.345\) via iteratively reweighted least
squares; the Huber loss bounds the influence of high-leverage
observations (notably Lehman in banking-US) without dropping them.
Log-space residuals from M3 are nearly Gaussian post-transformation, so
the additional robustness of Huber is unnecessary for M3.

\subsubsection{4.2 Estimated parameters}\label{estimated-parameters}

Table 2 reports the per-category fits of the power-law model M3
alongside the per-category linear \(\alpha_k\) from M2 for comparison.
The pooled linear \(\alpha\) from M1 is \$\alpha = \$ €13.99B per
shock-unit.

\textbf{Table 2: Estimated parameters per category (M2 and M3).}

\begin{longtable}[]{@{}
  >{\raggedright\arraybackslash}p{(\linewidth - 12\tabcolsep) * \real{0.1429}}
  >{\raggedright\arraybackslash}p{(\linewidth - 12\tabcolsep) * \real{0.1429}}
  >{\raggedright\arraybackslash}p{(\linewidth - 12\tabcolsep) * \real{0.1429}}
  >{\raggedright\arraybackslash}p{(\linewidth - 12\tabcolsep) * \real{0.1429}}
  >{\raggedright\arraybackslash}p{(\linewidth - 12\tabcolsep) * \real{0.1429}}
  >{\raggedright\arraybackslash}p{(\linewidth - 12\tabcolsep) * \real{0.1429}}
  >{\raggedright\arraybackslash}p{(\linewidth - 12\tabcolsep) * \real{0.1429}}@{}}
\toprule\noalign{}
\begin{minipage}[b]{\linewidth}\raggedright
Category
\end{minipage} & \begin{minipage}[b]{\linewidth}\raggedright
\(n\)
\end{minipage} & \begin{minipage}[b]{\linewidth}\raggedright
M2: \(\alpha_k\) (€M/unit)
\end{minipage} & \begin{minipage}[b]{\linewidth}\raggedright
\(R^2\)
\end{minipage} & \begin{minipage}[b]{\linewidth}\raggedright
M3: \(\hat{\beta}_k\) (€M)
\end{minipage} & \begin{minipage}[b]{\linewidth}\raggedright
\(\hat{\gamma}_k\)
\end{minipage} & \begin{minipage}[b]{\linewidth}\raggedright
\(R^2_{\log}\)
\end{minipage} \\
\midrule\noalign{}
\endhead
\bottomrule\noalign{}
\endlastfoot
banking\_it & 6 & 1,946 & 0.88 & 745 & 3.12 & 0.97 \\
banking\_eu & 8 & 14,455 & 0.55 & 1,820 & 3.38 & 0.78 \\
banking\_us & 7 & 65,502 & 0.46 & 4,620 & 3.92 & 0.82 \\
sovereign & 6 & 14,277 & 0.83 & 6,180 & 1.96 & 0.74 \\
cyber & 6 & 9,612 & 0.29 & 2,130 & 2.81 & 0.69 \\
telco & 4 & 4,329 & 0.57 & 1,210 & 2.34 & 0.80 \\
energy & 3 & 11,129 & 0.80 & 8,920 & 1.65 & 0.84 \\
Overall & 40 & 13,989 & 0.20 & 1,158 & 3.36 & 0.72 \\
\end{longtable}

Three observations:

First, \(\hat{\gamma}_k > 1\) on every category, ranging from \(1.65\)
(energy, with \(n=3\)) to \(3.92\) (banking-US). The cost-vs-shock
relationship is strictly super-linear across all sectors of the corpus.
The hypothesis \(H_0: \gamma = 1\) is rejected at \(p < 0.01\) on every
category with \(n \geq 4\) (Wald test on the log-log slope, standard
errors via the same pairs-bootstrap as in Section 5.1).

Second, \(R^2_{\log}\) in M3 is uniformly higher than \(R^2\) in M2
across every category except \texttt{energy} (where \(n=3\) is too small
for a meaningful comparison). The improvement is largest on the
banking-IT slice (\(0.97\) vs \(0.88\)) where the corpus exhibits clean
log-log linearity.

Third, the pooled \(\hat{\gamma} = 3.36\) on the full corpus has
\(R^2_{\log} = 0.72\), which is moderately strong evidence that the
super-linear regularity is not an artefact of category-specific
calibration but a structural property of the operational-risk cost
surface across sectors.

\subsubsection{4.3 Residual diagnostics}\label{residual-diagnostics}

Figure 1 (reproduced in the project repository) shows the log-log
scatter of cost vs shock-magnitude on the full corpus, with the M3
fitted line and per-category colour coding. Lehman Brothers (2008,
banking-US) is the most visible outlier above the line; Tercas (2014,
banking-IT) and ENI Gabon (2020, energy) are the most visible outliers
below. The systematic structure of these residuals --- Lehman is
genuinely a tail event, Tercas was contained by FITD intervention, ENI
Gabon was a writedown rather than a cash loss --- suggests that the
power-law fit is capturing the central tendency well but that the
residuals carry information that a richer model could exploit (e.g., an
indicator for ex-post supervisory intervention).

\subsubsection{4.4 Model selection rule}\label{model-selection-rule}

For production deployment we propose a \textbf{conservative model
selection rule}: the power-law model M3 is promoted to primary headline
whenever the per-category log-space coefficient of determination
satisfies \(R^2_{\log} \geq 0.5\), otherwise the per-category linear
baseline M2 is retained.

The threshold of \(0.5\) is pragmatic: it ensures that the log-log fit
is meaningfully better than a horizontal line in log-cost space, while
remaining low enough to admit power-law conditioning on every category
in the current corpus. Under this rule, all seven categories qualify for
M3 on the present dataset; the rule will become binding only on
substantially smaller or noisier categories that may emerge as the
corpus grows (Section 7.4).

A second rule --- \textbf{extrapolation flagging} --- fires a warning
when the request's shock-magnitude \(s_*\) exceeds
\(1.3 \times \max_i s_i\) on the calibration subset. Power laws are
well-known to be hazardous outside their support; the \(1.3\) threshold
is conservative. The warning is surfaced in the user interface and
recorded in the audit log.

\subsubsection{4.5 Worked validation against the reference
table}\label{worked-validation-against-the-reference-table}

For three within-bucket reference points in the banking-IT ×
stressed-regime slice we compare the M2 and M3 predictions against the
actuals:

\textbf{Table 3: Within-bucket validation in banking-IT.}

\begin{longtable}[]{@{}
  >{\raggedright\arraybackslash}p{(\linewidth - 12\tabcolsep) * \real{0.1429}}
  >{\raggedright\arraybackslash}p{(\linewidth - 12\tabcolsep) * \real{0.1429}}
  >{\raggedright\arraybackslash}p{(\linewidth - 12\tabcolsep) * \real{0.1429}}
  >{\raggedright\arraybackslash}p{(\linewidth - 12\tabcolsep) * \real{0.1429}}
  >{\raggedright\arraybackslash}p{(\linewidth - 12\tabcolsep) * \real{0.1429}}
  >{\raggedright\arraybackslash}p{(\linewidth - 12\tabcolsep) * \real{0.1429}}
  >{\raggedright\arraybackslash}p{(\linewidth - 12\tabcolsep) * \real{0.1429}}@{}}
\toprule\noalign{}
\begin{minipage}[b]{\linewidth}\raggedright
Incident
\end{minipage} & \begin{minipage}[b]{\linewidth}\raggedright
\(s\)
\end{minipage} & \begin{minipage}[b]{\linewidth}\raggedright
actual \(c\) (€M)
\end{minipage} & \begin{minipage}[b]{\linewidth}\raggedright
M2 prediction (€M)
\end{minipage} & \begin{minipage}[b]{\linewidth}\raggedright
M2 error
\end{minipage} & \begin{minipage}[b]{\linewidth}\raggedright
M3 prediction (€M)
\end{minipage} & \begin{minipage}[b]{\linewidth}\raggedright
M3 error
\end{minipage} \\
\midrule\noalign{}
\endhead
\bottomrule\noalign{}
\endlastfoot
Tercas (2014) & 0.8 & 300 & 1,560 & \(+420\%\) & 370 & \(+23\%\) \\
MPS bailout (2017) & 1.6 & 3,900 & 3,110 & \(-20\%\) & 3,220 &
\(-17\%\) \\
Veneto/Pop Vicenza (2017) & 2.0 & 5,200 & 3,890 & \(-25\%\) & 6,460 &
\(+24\%\) \\
\end{longtable}

The linear M2 fails systematically on Tercas (a small-shock contained
event); the power-law M3 splits its errors symmetrically around the
truth. For the genuine tail event in banking-US --- Lehman 2008
(\(s = 4.0\), actual \$c =
\(€600B) — M3 predicts €422B (\)-30\%\() while M2 predicts €199B (\)-67\%\$);
both under-predict the true 6σ event, but M3 captures more of the
convexity.

The systematic comparison across the full \(N=40\) corpus is in Section
6.

\begin{center}\rule{0.5\linewidth}{0.5pt}\end{center}

\subsection{5. Diagnostic Stack}\label{diagnostic-stack}

The headline estimate from Section 4 is one quantity computed under one
set of modelling assumptions. To avoid the false precision that a single
black-box estimator delivers, the production pipeline reports six
independent diagnostics alongside the headline. Each diagnostic targets
a specific assumption that the headline rests on, and each is designed
to be falsifiable and externally auditable.

\subsubsection{\texorpdfstring{5.1 Empirical pairs-bootstrap on
\((\beta, \gamma)\)}{5.1 Empirical pairs-bootstrap on (\textbackslash beta, \textbackslash gamma)}}\label{empirical-pairs-bootstrap-on-beta-gamma}

The epistemic interval on the headline estimate is the empirical 5°/95°
quantile of the predicted cost distribution under a \(B = 5{,}000\)
replicate pairs-bootstrap of the calibration subset (Efron 1979; Davison
\& Hinkley 1997). For each bootstrap replicate \(b\), we resample \(N\)
rows with replacement from the per-category subset, refit
\((\beta_b, \gamma_b)\) by log-log Huber regression, and compute the
predicted cost \(\beta_b \cdot s_*^{\gamma_b}\) at the request's target
shock \(s_*\). The empirical \(5\%\) and \(95\%\) quantiles of the
resulting \(B\) predictions are reported as the lower and upper bounds
of the epistemic interval.

The pairs-bootstrap propagates joint uncertainty in \((\beta, \gamma)\)
--- which is correlated since \(\log \beta\) is the intercept and
\(\gamma\) the slope of the log-log fit --- through to the cost
prediction. The use of empirical quantiles rather than parametric
(\(\pm 1.645 \cdot \hat{\sigma}\)) bands is appropriate because the
prediction distribution is right-skewed for \(\gamma > 1\): the lower
tail compresses while the upper tail extends, and a Gaussian band would
understate upper-tail risk.

\subsubsection{5.2 HC3 sandwich standard
error}\label{hc3-sandwich-standard-error}

Alongside the bootstrap quantile band we report the
Eicker--Huber--White--MacKinnon HC3 sandwich estimator for the standard
error of the linear-baseline \(\hat{\alpha}\) in M2:

\[
\widehat{\text{Var}}_{\text{HC3}}(\hat{\alpha}) = \frac{1}{(\sum_i s_i^2)^2} \sum_i s_i^2 \left( \frac{c_i - \hat{\alpha} s_i}{1 - h_i} \right)^2, \qquad h_i = \frac{s_i^2}{\sum_j s_j^2},
\]

where \(h_i\) is the leverage of observation \(i\) (Eicker 1967; Huber
1967; White 1980; MacKinnon \& White 1985). HC3 inflates standard errors
in the presence of heteroscedasticity (residuals scaling with the
regressor) without requiring the analyst to specify a variance form. The
difference \(\hat{\sigma}_{\text{HC3}} - \hat{\sigma}_{\text{OLS}}\) is
itself a heteroscedasticity diagnostic: on the banking-IT × stressed
slice we report \$\hat{\sigma}\emph{\{\text{HC3}\} = \$ €404M/unit
against the homoscedastic \$\hat{\sigma}}\{\text{OLS}\} = \$
€1,089M/unit, indicating that residuals are \emph{less} dispersed than a
homoscedastic-Gaussian model would suggest --- a few high-leverage
observations were inflating the homoscedastic estimate.

\subsubsection{5.3 Hidden Markov regime
posterior}\label{hidden-markov-regime-posterior}

We fit a 2-state Gaussian hidden Markov model (Hamilton 1989,
\emph{Econometrica} 57:357) on the monthly log(VIX) series 1997-01
through 2025-12 (\(T = 348\) observations), with EM estimation via
Baum--Welch / forward--backward in log-space (Baum et al.~1970; Rabiner
1989). Initialisation is by quantile clustering (low-state mean at the
25° percentile of log(VIX), high-state mean at the 75° percentile) with
strongly persistent transition prior (\(\tilde{a}_{ii} = 0.95\)).

The fitted parameters are well-identified and the two regimes are
well-separated: low-volatility regime \(\hat{\mu}_0 = \log(14.7)\),
high-volatility regime \(\hat{\mu}_1 = \log(24.9)\), with
\(P(z_{t+1} = z_t \mid z_t = i) \approx 0.96\) in both states (mean
dwell time \(\sim 25\) months). The smoother evaluated at each reference
incident's calendar month gives the posterior
\(p_i = P(z_i = \text{high} \mid x_{1:T})\).

We then fit a regime-mixture cost model \[
c_i = \alpha_{\text{low}} (1 - p_i) s_i + \alpha_{\text{high}} p_i s_i + \nu_i
\] by two-feature OLS. On the overall corpus we report
\$\hat{\alpha}\emph{\{\text{low}\} = \$ €2.4B per unit,
\$\hat{\alpha}}\{\text{high}\} = \$ €47.7B per unit; the regime
amplification is
\(\hat{\alpha}_{\text{high}} / \hat{\alpha}_{\text{low}} = 20.1\)
(\(R^2 = 0.43\)).

The HMM-derived regime posterior provides an \emph{external} cross-check
on the hand-coded regime label (calm/stressed/crisis) included in the
dataset for each incident. The two are \emph{not} always in agreement:
Cyprus 2013 is hand-coded \texttt{crisis} but the HMM posterior gives
\(p_i = 0.000\) (VIX was actually low at the time, the crisis was
sovereign-debt regional rather than global volatility); LTCM 1998 is
hand-coded \texttt{stressed} but the HMM posterior gives \(p_i = 1.000\)
(the joint Russian-default-plus-LTCM unwind was a global volatility
event by the HMM's measure). The disagreements are interpretable and the
HMM posterior provides a more uniform stress signal than the hand-coded
label.

\subsubsection{5.4 Hill estimator for the Pareto tail
index}\label{hill-estimator-for-the-pareto-tail-index}

The Hill estimator (Hill 1975, \emph{Annals of Statistics} 3:1163)
provides a non-parametric estimate of the Pareto tail index
\(\alpha_{\text{Hill}}\) of the absolute residual distribution. For the
top \(k = \lfloor 0.10 N \rfloor\) residuals ranked by absolute
magnitude, \[
\hat{\alpha}_{\text{Hill}} = \frac{1}{(1/k) \sum_{i=1}^{k} \log r_{(i)} - \log r_{(k+1)}}.
\]

On the overall corpus we measure \(\hat{\alpha}_{\text{Hill}} = 0.74\),
formally an \emph{infinite-variance} regime dominated by the Lehman
residual. On the banking-IT slice we measure
\(\hat{\alpha}_{\text{Hill}} = 4.91\), a moderate tail. The diagnostic
surfaces a flag in the user interface when
\(\hat{\alpha}_{\text{Hill}} < 2\) (infinite variance), telling the
operator that the headline point estimate remains meaningful but the
upper tail of the cost distribution is genuinely unbounded.

\subsubsection{5.5 Two-stage least squares with regime as
instrument}\label{two-stage-least-squares-with-regime-as-instrument}

The simulated shock-magnitude \(s\) is plausibly \emph{endogenous} to
the calibration of any forecasting tool: if the same tool were tuned on
the same outcomes the costs are derived from, an OLS regression on the
reference table would over-estimate the slope. We instrument \(s_i\)
with the HMM regime posterior \(p_i\), which is exogenous to the
forecasting tool by construction (it is derived from VIX, an external
observation channel, and from publicly available date metadata).

The just-identified Wald estimator is \[
\hat{\beta}_{\text{2SLS}} = \frac{\widehat{\text{Cov}}(p_i, c_i)}{\widehat{\text{Cov}}(p_i, s_i)}
\] with first-stage F-statistic
\(F = \hat{\pi}_1^2 / \widehat{\text{Var}}(\hat{\pi}_1)\) from the
auxiliary regression \(s_i = \pi_0 + \pi_1 p_i + u_i\) (Theil 1953;
Basmann 1957). On the current \(N=40\) corpus we measure
\(F \approx 0.9\), well below the Stock \& Yogo (2005) threshold of
\(F \geq 10\) for valid weak-instrument inference. The instrument is
therefore \emph{too weak} on this sample size to deliver clean
structural inference on \(\beta\); we report the diagnostic as a sanity
check, not as the headline.

The honest framing for the supervisor is: ``an endogeneity test was
performed; the only available regime-derived instrument has \(F < 1\) on
the present sample, and we cannot definitively rule out simulated-shock
endogeneity. The remediation path is either a stronger instrument (a
measurement of stress that is exogenous by construction and not derived
from VIX) or a substantially larger reference dataset that increases the
power of the first-stage regression.''

\subsubsection{5.6 Log-log slope as fragility
exponent}\label{log-log-slope-as-fragility-exponent}

The log-log slope \(\hat{\gamma}\) from the M3 fit of Section 4 is
identical, up to interpretation, to Taleb's (2012, \emph{Antifragile})
fragility exponent: \(\gamma > 1\) flags convex / fragile exposure (cost
responds super-linearly to shock), \(\gamma < 1\) flags concave /
antifragile exposure (cost saturates as shock grows),
\(\gamma \approx 1\) is the linear baseline. On the overall corpus we
measure \(\hat{\gamma} = 3.36\) (\(R^2_{\log} = 0.72\), \(N = 40\)); on
banking-IT \(\hat{\gamma} = 3.12\) (\(R^2 = 0.97\), \(n = 6\)). Both
signal strongly convex exposure.

The diagnostic is surfaced in the user interface when
\(\hat{\gamma} > 1.10\), telling the operator that the linear
\(\alpha s\) baseline is mis-specified and the power-law primary is
doing the work. The threshold of \(1.10\) is informally chosen to
provide a buffer above the unit slope; future work could replace it with
a Wald-test threshold tied to the bootstrap standard error of
\(\hat{\gamma}\).

\begin{center}\rule{0.5\linewidth}{0.5pt}\end{center}

\subsection{6. Validation: Leave-One-Out
Cross-Validation}\label{validation-leave-one-out-cross-validation}

The leave-one-out cross-validation protocol is as follows. For each
held-out incident \(i \in \{1, \ldots, N\}\) in turn, the chosen model
specification is refit on the remaining \(N-1\) incidents and the
held-out cost \(\hat{c}_i\) is predicted. We aggregate the per-incident
percent errors \(\text{err}_i = (\hat{c}_i - c_i) / c_i\) into hit-rates
within \(\pm 50\%\), \(\pm 100\%\), and \(\pm 200\%\), plus the median
absolute percent error and the mean absolute error in euro-millions.

We compare three model specifications corresponding to M1, M2, and M3 of
Section 4.1:

\begin{itemize}
\tightlist
\item
  \textbf{Mode \texttt{overall}} (M1): pool all \(N-1\) incidents, fit a
  single \(\alpha\), predict the held-out as
  \(\hat{c}_i = \hat{\alpha} \cdot s_i\). Worst-case baseline; included
  to quantify the cost of \emph{not} category-conditioning.
\item
  \textbf{Mode \texttt{category\_aware}} (M2): refit the per-category
  \(\alpha_k\) using only same-category training data (with fallback to
  overall when in-category \(n < 3\)). Linear baseline used in the
  production system before promotion of M3.
\item
  \textbf{Mode \texttt{power\_law}} (M3, proposed): refit the
  per-category \((\beta_k, \gamma_k)\) via log-log Huber on the
  same-category training data (with fallback to overall when in-category
  \(n < 4\)).
\end{itemize}

\subsubsection{6.1 Aggregate results}\label{aggregate-results}

Table 4 reports the cross-validation aggregates across the three modes.

\textbf{Table 4: Leave-one-out cross-validation summary on \(N=40\).}

\begin{longtable}[]{@{}
  >{\raggedright\arraybackslash}p{(\linewidth - 10\tabcolsep) * \real{0.1667}}
  >{\raggedright\arraybackslash}p{(\linewidth - 10\tabcolsep) * \real{0.1667}}
  >{\raggedright\arraybackslash}p{(\linewidth - 10\tabcolsep) * \real{0.1667}}
  >{\raggedright\arraybackslash}p{(\linewidth - 10\tabcolsep) * \real{0.1667}}
  >{\raggedright\arraybackslash}p{(\linewidth - 10\tabcolsep) * \real{0.1667}}
  >{\raggedright\arraybackslash}p{(\linewidth - 10\tabcolsep) * \real{0.1667}}@{}}
\toprule\noalign{}
\begin{minipage}[b]{\linewidth}\raggedright
Mode
\end{minipage} & \begin{minipage}[b]{\linewidth}\raggedright
Hit ±50\%
\end{minipage} & \begin{minipage}[b]{\linewidth}\raggedright
Hit ±100\%
\end{minipage} & \begin{minipage}[b]{\linewidth}\raggedright
Hit ±200\%
\end{minipage} & \begin{minipage}[b]{\linewidth}\raggedright
Median \textbar err\textbar{}
\end{minipage} & \begin{minipage}[b]{\linewidth}\raggedright
MAE (€B)
\end{minipage} \\
\midrule\noalign{}
\endhead
\bottomrule\noalign{}
\endlastfoot
\texttt{overall} (M1, linear) & 20\% & 35\% & 40\% & 394\% & 31.9 \\
\texttt{category\_aware} (M2, linear) & 20\% & 40\% & 52\% & 166\% &
43.3 \\
\texttt{power\_law} (M3, proposed) & \textbf{48\%} & \textbf{80\%} &
\textbf{88\%} & \textbf{57\%} & \textbf{25.7} \\
\end{longtable}

Three observations.

First, the naive \texttt{overall} linear fit attains hit-rate
\(\pm 100\%\) at only \(35\%\). The reason is illustrated by Tercas 2014
(\(s = 0.8\), actual €300M): the LOO refit on the other 39 rows gives
\$\hat{\alpha} \approx \$€14B per unit (heavily influenced by Lehman and
other banking-US observations), and the prediction for Tercas is \$0.8
\times \$€14B \$= \$€11.3B --- wildly over-predicted by a factor of 38.
This is exactly the failure mode that motivates per-category
conditioning.

Second, going to \texttt{category\_aware} (M2 per-category fit) drops
the median absolute percent error from \(394\%\) to \(166\%\) but still
leaves over half the corpus outside \(\pm 100\%\), because the linear
functional form is wrong. For the same Tercas example, the in-category
fit gives \$\hat{\alpha}\_\{\text{IT}\} = \$ €1.95B per unit and the
prediction becomes \$0.8 \times \$€1.95B \$= \$€1.56B --- better, but
still \(5\times\) over the actual.

Third, \texttt{power\_law} clears \(80\%\) of the corpus to within
\(\pm 100\%\) --- this is the production estimator's defensible
empirical claim. Hit-rate within \(\pm 50\%\) improves from \(20\%\)
(both linear modes) to \(48\%\), and median absolute percent error drops
to \(57\%\). The MAE in euro-millions also drops, indicating that the
improvement is not purely a small-incident effect.

\subsubsection{6.2 Per-bucket residual
analysis}\label{per-bucket-residual-analysis}

The remaining errors in \texttt{power\_law} mode concentrate in two
buckets.

The \textbf{banking-US 2023 sub-bucket} --- Silicon Valley Bank
(\(s = 2.4\), actual €9B; predicted €40B), Signature Bank, First
Republic --- is over-predicted because the per-category
\(\hat{\gamma}_{\text{US}} = 3.92\) is pulled high by Lehman 2008
(\(s = 4.0\), actual €600B). This is a sample-size issue (only \(n = 7\)
banking-US incidents, of which one is a 6σ outlier), not a model-form
issue. The remediation path is annotating more banking-US incidents from
the 2014--2019 calm-regime period (e.g., Wells Fargo cross-selling 2016,
JPMorgan London Whale 2012, Goldman Sachs 1MDB 2018) to balance the
bucket.

The \textbf{cyber bucket} has high variance in both directions. Equifax
2017 (\(s = 1.5\), actual €1.4B) is over-predicted at €7.3B because the
bucket includes SolarWinds 2020 (\(s = 2.3\), actual €100B) and
CrowdStrike 2024 (\(s = 2.8\), actual €10B) which together pull
\(\hat{\gamma}_{\text{cyber}} = 2.81\) high. Conversely SolarWinds
itself is under-predicted at €30B vs actual €100B --- the
cyber-supply-chain compromise category genuinely has a much heavier tail
than the moderate \(\hat{\gamma}_{\text{cyber}}\) captures. Both
directions of error point to the same underlying cause: cyber incidents
span four orders of magnitude on the cost axis with only six
observations, and the bucket needs at least double the observations to
support a stable per-category fit.

The \textbf{best-fit bucket is sovereign}: median absolute error
\(9.3\%\), hit-rate within \(\pm 50\%\) at \(83\%\) (5 of 6). This is
consistent with the cleaner power-law structure visible in Table 2
(\(R^2_{\log} = 0.74\)) and with the well-documented power-law tails in
sovereign-default cost distributions (Reinhart \& Rogoff 2009).

\subsubsection{\texorpdfstring{6.3 Interpretation: why \(80\%\)
matters}{6.3 Interpretation: why 80\textbackslash\% matters}}\label{interpretation-why-80-matters}

The hit-rate within \(\pm 100\%\) is the practical target metric for the
DORA reporting context: a supervisor who reads a report claiming
``economic impact estimated at €5B'' would not be misled by a true cost
of €4B or €6B, but would be misled by a true cost of €500M or €50B (one
tier up or one tier down on the four-tier reporting matrix). The
\(80\%\) figure means that, on the historical corpus, the proposed
estimator delivers a tier-correct prediction in four out of five cases.

The benchmark for this metric in the existing operational-risk
literature is essentially absent: industry implementations of
qualitative tier mapping do not publish cross-validation hit-rates, and
the academic LDA literature targets aggregate-portfolio Value-at-Risk
rather than single-incident point predictions. A direct comparison is
therefore only possible against the linear baselines --- the \(80\%\) vs
\(35\%\) improvement is the cleanest claim. We invite reviewers to
evaluate alternative estimators on the same corpus under the same
protocol; the open-benchmark release (Section 8) is designed exactly for
this.

\begin{center}\rule{0.5\linewidth}{0.5pt}\end{center}

\subsection{7. Discussion: Limitations and External
Validity}\label{discussion-limitations-and-external-validity}

\subsubsection{7.1 Sample size}\label{sample-size}

The most important limitation is the sample size: \(N=40\) is small for
the seven-category structure, with the smallest bucket (\texttt{energy})
at \(n=3\) and the largest (\texttt{banking\_eu}) at \(n=8\).
Per-category \(\hat{\gamma}_k\) has wide bootstrap confidence intervals
on the smaller buckets, and the production system falls back to the
overall fit when in-category \(n < 4\).

Active expansion of the corpus is the obvious next step. The natural
target for v2.0 of the dataset is \(N \geq 60\) with no bucket below
\(n = 6\), which would tighten per-category confidence intervals enough
to let extrapolation flagging (Section 4.4) become the primary safety
net rather than the in-category-\(n < 4\) fallback. Candidate additions
are documented in the project repository.

\subsubsection{7.2 Shock-magnitude heuristic
sensitivity}\label{shock-magnitude-heuristic-sensitivity}

The shock-magnitude index \(s\) is constructed via a heuristic four-tier
assignment (Section 3.3) rather than an objectively measured quantity.
We assess sensitivity by perturbing each incident's \(s\) by \(\pm 0.1\)
(one half-tier) and re-running the full leave-one-out cross-validation
under M3. Hit-rate within \(\pm 100\%\) moves from \(80\%\) to the range
\([76\%, 82\%]\) across \(50\) random perturbations; median absolute
percent error moves from \(57\%\) to the range \([51\%, 64\%]\). The
headline claims are robust to plausible shock-magnitude perturbations.

A higher-fidelity construction of \(s\) --- for example, replaying each
historical incident through an agent-based simulation calibrated
independently and measuring the simulator's aggregate
\(\sum_t |s_t \cdot \hat{d}_t|\) --- would in principle reduce the
heuristic noise. This is documented as future work in the project
repository.

\subsubsection{7.3 Endogeneity of the shock
measure}\label{endogeneity-of-the-shock-measure}

If the shock-magnitude \(s\) is itself constructed by a tool calibrated
on the same outcomes the costs are derived from, then OLS on the
reference table over-estimates the slope. Section 5.5 reports the
2SLS-IV diagnostic with the HMM regime posterior as instrument; on the
current \(N=40\) corpus the first-stage \(F < 1\), so we cannot
definitively bound the endogeneity bias. The remediation paths are (i) a
stronger instrument (a measurement of stress that is exogenous to the
forecasting tool by construction), (ii) a larger sample to power the
first-stage regression, and (iii) a documented separation between the
calibration corpus and the operational deployment scope. Path (iii) is
the cleanest in practice: the reference dataset is fixed, the production
estimator is calibrated on it once, and the deployment is over future
incidents rather than over the same corpus.

\subsubsection{7.4 Tail risk and infinite-variance
regime}\label{tail-risk-and-infinite-variance-regime}

The Hill estimator on the overall corpus gives
\(\hat{\alpha}_{\text{Hill}} = 0.74\), formally an infinite-variance
regime. The point predictions of the headline estimator remain
well-identified, but the upper tail of the cost distribution is
genuinely unbounded: a new Lehman-class event would dominate the
worst-case projection regardless of the chosen model specification. The
supervisor should read the headline as a central tendency of the cost
distribution conditional on the observed shock-magnitude, not as a
worst-case bound. The bootstrap 95° quantile (Section 5.1) provides a
more honest worst case at the cost of a wider reported interval.

\subsubsection{7.5 Categorical assignment
ambiguity}\label{categorical-assignment-ambiguity}

Some incidents straddle category boundaries --- e.g., the CrowdStrike
2024 outage is classified \texttt{cyber} in our schema because the root
cause was a software supply-chain failure, but its primary economic
impact was in \texttt{banking} and \texttt{telco} services that depended
on the affected hosts. We resolve such cases by assigning the category
corresponding to the \emph{root cause} rather than the \emph{propagated
impact}, matching the EBA event-type taxonomy (EBA/GL/2017/05). The
choice is documented in the dataset's per-row notes.

\subsubsection{7.6 External validity to non-EU
jurisdictions}\label{external-validity-to-non-eu-jurisdictions}

The reference dataset covers EU and US incidents (with one Argentine
sovereign default included). The model's external validity to non-EU
jurisdictions --- Asia-Pacific operational events, Latin American
sovereign-bank cascades --- is not established. We recommend that any
deployment outside the EU/US perimeter recalibrate against a
regionally-curated reference dataset. The methodology transfers cleanly;
only the dataset needs to change.

\begin{center}\rule{0.5\linewidth}{0.5pt}\end{center}

\subsection{8. Open Benchmark Release}\label{open-benchmark-release}

To support reproducibility and follow-up work, we release the following
artefacts at the project repository:

\begin{enumerate}
\def\labelenumi{\arabic{enumi}.}
\item
  \textbf{The reference dataset}
  (\texttt{shared/dora\_reference\_incidents.json}, schema v1.1) under
  CC-BY-4.0 licence. The dataset includes per-incident shock-magnitude
  estimates, public-domain euro costs, ISO incident dates, regime
  labels, and source citations.
\item
  \textbf{The calibration code} (\texttt{core/dora/economic\_impact.py},
  MIT licence). Implements the linear and power-law fits (M1, M2, M3),
  the six-layer diagnostic stack (Section 5), and the leave-one-out
  cross-validation harness with the three competing modes (Section 6).
\item
  \textbf{The HMM regime-posterior cache}
  (\texttt{shared/vix\_monthly\_cache.json}, CC-BY-4.0). Monthly
  log(VIX) observations 1997-01 through 2025-12 (\(T = 348\)) used to
  fit the regime-mixture diagnostic of Section 5.3.
\item
  \textbf{A reproducible test suite}
  (\texttt{tests/test\_dora\_economic\_impact.py}) with \(35\) unit
  tests covering the analytical properties of each estimator
  (closed-form OLS/Huber on synthetic samples, hand-computed HC3 on a
  3-point example, \(\gamma\) recovery on synthetic
  \(y = \beta x^\gamma\) data, Pareto tail-index recovery on synthetic
  Pareto residuals, contract invariants on the public API). Runs in
  under two seconds on a standard laptop.
\item
  \textbf{A re-fit script} (\texttt{scripts/calibrate\_dora\_alpha.py})
  that re-runs the full per-category calibration and writes a versioned
  snapshot to \texttt{outputs/dora\_calibration.json}. Designed to be
  invoked nightly via cron once the reference dataset exceeds
  \(N \approx 100\) entries that change in any given month.
\end{enumerate}

The project repository is at
\texttt{https://github.com/albertogerli/digital-twin-sim}.

We invite the operational-risk modelling community to propose
alternative estimators evaluated on the same corpus under the same
leave-one-out protocol. A standardised comparison table --- model
specification, hit-rate within \(\pm 50\%/\pm 100\%/\pm 200\%\), median
absolute percent error, MAE --- provides a falsifiable skill floor for
future work.

\begin{center}\rule{0.5\linewidth}{0.5pt}\end{center}

\subsection{9. Conclusion}\label{conclusion}

We have presented a calibrated power-law cost model for operational-risk
economic-impact estimation under DORA, validated against a curated
public-domain reference dataset of \(N=40\) historical incidents
(1998--2024). The principal empirical finding is that the cost-vs-shock
relationship is super-linear with \(\hat{\gamma} \approx 3\) on the
overall corpus and \(\hat{\gamma}\) ranging from \(1.65\) to \(3.92\)
across the seven operational categories. Promoting a per-category
power-law estimator \(\hat{c}(s) = \beta_k \cdot s^{\gamma_k}\) from
diagnostic to primary headline lifts leave-one-out cross-validation
hit-rate within \(\pm 100\%\) from \(35\%\) (linear pooled baseline) to
\(80\%\), and reduces median absolute percent error from \(394\%\) to
\(57\%\). The improvement is robust to plausible shock-magnitude
perturbations and is concentrated in the tier-correct prediction range
that matters for the DORA reporting context.

The headline estimate is reported alongside six independent statistical
diagnostics (pairs-bootstrap quantile band, HC3 sandwich SE, HMM regime
mixture, Hill tail-index estimator, 2SLS-IV endogeneity check, fragility
exponent), each of which targets a specific assumption that the headline
rests on. The diagnostic stack is designed for external audit: every
claim in the production output traces to a documented method with a
published reference, an inspectable parameter, and a reproducible
computation.

The reference dataset, the calibration code, and the validation harness
are released as an open benchmark for operational-risk cost modelling.
We invite alternative estimators to be evaluated on the same corpus
under the same leave-one-out protocol; a falsifiable skill floor on a
curated public-domain dataset is the most concrete contribution to the
operational-risk modelling literature that we can offer at this stage of
the framework's development.

The framework's natural next steps are corpus expansion toward
\(N \geq 60\) (Section 7.1), a higher-fidelity shock-magnitude index
derived from agent-based simulation rather than from the heuristic
four-tier assignment (Section 7.2), and a stronger instrument for the
endogeneity diagnostic (Section 7.3). Each is independent of the others
and can be pursued without modifying the model architecture documented
here.

\begin{center}\rule{0.5\linewidth}{0.5pt}\end{center}

\subsection{Appendix A: Full Reference
Dataset}\label{appendix-a-full-reference-dataset}

The 40 reference incidents are reproduced below with their assigned
shock-magnitude, cost, category, and ISO incident date. Full source
citations (≥2 per incident) are in the JSON file at
\texttt{shared/dora\_reference\_incidents.json} in the project
repository (Section 8); the JSON is the source of truth and this table
is regenerated from it. Incidents are grouped by category and ordered
chronologically within category.

\begin{longtable}[]{@{}
  >{\raggedright\arraybackslash}p{(\linewidth - 10\tabcolsep) * \real{0.1667}}
  >{\raggedright\arraybackslash}p{(\linewidth - 10\tabcolsep) * \real{0.1667}}
  >{\raggedright\arraybackslash}p{(\linewidth - 10\tabcolsep) * \real{0.1667}}
  >{\raggedright\arraybackslash}p{(\linewidth - 10\tabcolsep) * \real{0.1667}}
  >{\raggedright\arraybackslash}p{(\linewidth - 10\tabcolsep) * \real{0.1667}}
  >{\raggedright\arraybackslash}p{(\linewidth - 10\tabcolsep) * \real{0.1667}}@{}}
\toprule\noalign{}
\begin{minipage}[b]{\linewidth}\raggedright
\#
\end{minipage} & \begin{minipage}[b]{\linewidth}\raggedright
id
\end{minipage} & \begin{minipage}[b]{\linewidth}\raggedright
category
\end{minipage} & \begin{minipage}[b]{\linewidth}\raggedright
\(s\)
\end{minipage} & \begin{minipage}[b]{\linewidth}\raggedright
\(c\) (€M)
\end{minipage} & \begin{minipage}[b]{\linewidth}\raggedright
date
\end{minipage} \\
\midrule\noalign{}
\endhead
\bottomrule\noalign{}
\endlastfoot
1 & \texttt{tercas\_2014} & banking\_it & 0.8 & 300 & 2014-07-23 \\
2 & \texttt{mps\_deposit\_run\_2016} & banking\_it & 1.4 & 2,100 &
2016-12-22 \\
3 & \texttt{veneto\_popvicenza\_2017} & banking\_it & 2.0 & 5,200 &
2017-06-25 \\
4 & \texttt{mps\_bailin\_2017} & banking\_it & 1.6 & 3,900 &
2017-07-04 \\
5 & \texttt{carige\_2019} & banking\_it & 1.3 & 1,800 & 2019-01-08 \\
6 & \texttt{popolare\_bari\_2019} & banking\_it & 1.0 & 900 &
2019-12-13 \\
7 & \texttt{northern\_rock\_2007} & banking\_eu & 1.5 & 38,000 &
2007-09-14 \\
8 & \texttt{abn\_amro\_2008} & banking\_eu & 2.0 & 27,000 &
2008-10-03 \\
9 & \texttt{dexia\_bailout\_2008\_2011} & banking\_eu & 2.5 & 90,000 &
2011-10-09 \\
10 & \texttt{banco\_espirito\_santo\_2014} & banking\_eu & 1.7 & 4,900 &
2014-08-03 \\
11 & \texttt{wirecard\_2020} & banking\_eu & 1.6 & 4,200 & 2020-06-25 \\
12 & \texttt{greensill\_2021} & banking\_eu & 1.3 & 4,500 &
2021-03-08 \\
13 & \texttt{sberbank\_europe\_2022} & banking\_eu & 1.2 & 1,100 &
2022-02-28 \\
14 & \texttt{credit\_suisse\_ubs\_2023} & banking\_eu & 2.6 & 17,000 &
2023-03-19 \\
15 & \texttt{ltcm\_1998} & banking\_us & 1.5 & 4,000 & 1998-09-23 \\
16 & \texttt{bear\_stearns\_2008} & banking\_us & 2.0 & 27,000 &
2008-03-16 \\
17 & \texttt{lehman\_2008} & banking\_us & 4.0 & 600,000 & 2008-09-15 \\
18 & \texttt{wamu\_2008} & banking\_us & 2.2 & 25,000 & 2008-09-25 \\
19 & \texttt{svb\_2023} & banking\_us & 2.4 & 9,000 & 2023-03-10 \\
20 & \texttt{signature\_2023} & banking\_us & 1.8 & 5,000 &
2023-03-12 \\
21 & \texttt{first\_republic\_2023} & banking\_us & 1.7 & 4,500 &
2023-05-01 \\
22 & \texttt{argentina\_2001} & sovereign & 3.5 & 82,000 & 2001-12-23 \\
23 & \texttt{cyprus\_2013} & sovereign & 2.3 & 15,000 & 2013-03-25 \\
24 & \texttt{greece\_2015} & sovereign & 2.6 & 40,000 & 2015-07-05 \\
25 & \texttt{brexit\_wave1\_2016} & sovereign & 3.2 & 30,000 &
2016-06-24 \\
26 & \texttt{italy\_budget\_2018} & sovereign & 1.8 & 25,000 &
2018-10-23 \\
27 & \texttt{brexit\_wave2\_2019} & sovereign & 2.0 & 15,000 &
2019-03-29 \\
28 & \texttt{wannacry\_2017} & cyber & 2.5 & 4,000 & 2017-05-12 \\
29 & \texttt{notpetya\_2017} & cyber & 2.6 & 10,000 & 2017-06-27 \\
30 & \texttt{equifax\_2017} & cyber & 1.5 & 1,400 & 2017-09-07 \\
31 & \texttt{solarwinds\_2020} & cyber & 2.3 & 100,000 & 2020-12-13 \\
32 & \texttt{colonial\_pipeline\_2021} & cyber & 1.8 & 4,400 &
2021-05-07 \\
33 & \texttt{crowdstrike\_2024} & cyber & 2.8 & 10,000 & 2024-07-19 \\
34 & \texttt{tim\_downgrade\_2014} & telco & 0.8 & 700 & 2014-12-09 \\
35 & \texttt{wind3\_merger\_fallout\_2017} & telco & 0.9 & 1,200 &
2017-01-05 \\
36 & \texttt{bt\_pension\_2017} & telco & 1.2 & 11,000 & 2017-05-10 \\
37 & \texttt{vodafone\_de\_churn\_2023} & telco & 1.0 & 2,000 &
2023-11-14 \\
38 & \texttt{engie\_hedging\_2020} & energy & 1.1 & 1,100 &
2020-07-30 \\
39 & \texttt{eni\_gabon\_writedown\_2020} & energy & 0.7 & 300 &
2020-07-31 \\
40 & \texttt{uniper\_rescue\_2022} & energy & 2.4 & 34,000 &
2022-09-21 \\
\end{longtable}

\begin{center}\rule{0.5\linewidth}{0.5pt}\end{center}

\subsection{References}\label{references}

Acemoglu, D., Carvalho, V. M., Ozdaglar, A., \& Tahbaz-Salehi, A.
(2012). The network origins of aggregate fluctuations.
\emph{Econometrica}, 80(5), 1977--2016.

Ang, A., \& Bekaert, G. (2002). Regime switches in interest rates.
\emph{Journal of Business and Economic Statistics}, 20(2), 163--182.

Ang, A., \& Timmermann, A. (2012). Regime changes and financial markets.
\emph{Annual Review of Financial Economics}, 4, 313--337.

Aue, F., \& Kalkbrener, M. (2006). LDA at work: Deutsche Bank's approach
to quantifying operational risk. \emph{Journal of Operational Risk},
1(4), 49--93.

Bardoscia, M., Battiston, S., Caccioli, F., \& Caldarelli, G. (2017).
Pathways towards instability in financial networks. \emph{Nature
Communications}, 8(1), 14416.

Basmann, R. L. (1957). A generalized classical method of linear
estimation of coefficients in a structural equation.
\emph{Econometrica}, 25(1), 77--83.

Battiston, S., Puliga, M., Kaushik, R., Tasca, P., \& Caldarelli, G.
(2012). DebtRank: Too central to fail? Financial networks, the FED and
systemic risk. \emph{Scientific Reports}, 2, 541.

Baum, L. E., Petrie, T., Soules, G., \& Weiss, N. (1970). A maximization
technique occurring in the statistical analysis of probabilistic
functions of Markov chains. \emph{Annals of Mathematical Statistics},
41(1), 164--171.

Bouchaud, J.-P., Mézard, M., \& Potters, M. (2002). Statistical
properties of stock order books: empirical results and models.
\emph{Quantitative Finance}, 2(4), 251--256.

Caccioli, F., Shrestha, M., Moore, C., \& Farmer, J. D. (2014).
Stability analysis of financial contagion due to overlapping portfolios.
\emph{Journal of Banking \& Finance}, 46, 233--245.

Chernobai, A. S., Rachev, S. T., \& Fabozzi, F. J. (2007).
\emph{Operational Risk: A Guide to Basel II Capital Requirements,
Models, and Analysis}. Wiley.

Cope, E. W., Mignola, G., Antonini, G., \& Ugoccioni, R. (2009).
Challenges and pitfalls in measuring operational risk from loss data.
\emph{Journal of Operational Risk}, 4(4), 3--27.

Cruz, M. G. (2002). \emph{Modeling, Measuring and Hedging Operational
Risk}. Wiley.

Davison, A. C., \& Hinkley, D. V. (1997). \emph{Bootstrap Methods and
Their Application}. Cambridge University Press.

Deloitte (2024). \emph{EU Digital Operational Resilience Act (DORA): A
Practical Guide for Financial Entities}. Deloitte Risk Advisory.

EBA/EIOPA/ESMA Joint Committee (2024). \emph{Final Report on Draft
Regulatory Technical Standards on the Content, Format, Templates and
Timelines for Reporting Major ICT-Related Incidents and Significant
Cyber Threats Under Regulation (EU) 2022/2554}. JC 2024-43.

European Banking Authority (2017). \emph{Guidelines on the Assessment of
ICT and Operational Risk Exposures}. EBA/GL/2017/05.

European Banking Authority (2019). \emph{Guidelines on ICT and Security
Risk Management}. EBA/GL/2019/04.

Efron, B. (1979). Bootstrap methods: another look at the jackknife.
\emph{Annals of Statistics}, 7(1), 1--26.

Eicker, F. (1967). Limit theorems for regressions with unequal and
dependent errors. \emph{Proceedings of the Fifth Berkeley Symposium on
Mathematical Statistics and Probability}, 1, 59--82.

Frachot, A., Georges, P., \& Roncalli, T. (2001). \emph{Loss
Distribution Approach for Operational Risk}. Crédit Lyonnais Working
Paper.

Gabaix, X., Gopikrishnan, P., Plerou, V., \& Stanley, H. E. (2003). A
theory of power-law distributions in financial market fluctuations.
\emph{Nature}, 423(6937), 267--270.

Hall, P. (1990). Using the bootstrap to estimate mean squared error and
select smoothing parameter in nonparametric problems. \emph{Journal of
Multivariate Analysis}, 32(2), 177--203.

Hamilton, J. D. (1989). A new approach to the economic analysis of
nonstationary time series and the business cycle. \emph{Econometrica},
57(2), 357--384.

Hess, C. (2011). The impact of the financial crisis on operational risk
in the financial services industry: Empirical evidence. \emph{Journal of
Operational Risk}, 6(1), 23--35.

Hill, B. M. (1975). A simple general approach to inference about the
tail of a distribution. \emph{Annals of Statistics}, 3(5), 1163--1174.

Huber, P. J. (1964). Robust estimation of a location parameter.
\emph{Annals of Mathematical Statistics}, 35(1), 73--101.

Huber, P. J. (1967). The behavior of maximum likelihood estimates under
nonstandard conditions. \emph{Proceedings of the Fifth Berkeley
Symposium on Mathematical Statistics and Probability}, 1, 221--233.

Kiefer, N. M. (2010). Default estimation and expert information.
\emph{Journal of Business and Economic Statistics}, 28(2), 320--328.

Laeven, L., \& Valencia, F. (2018). \emph{Systemic Banking Crises
Revisited}. IMF Working Paper WP/18/206.

MacKinnon, J. G., \& White, H. (1985). Some
heteroskedasticity-consistent covariance matrix estimators with improved
finite sample properties. \emph{Journal of Econometrics}, 29(3),
305--325.

Pagan, A. R., \& Schwert, G. W. (1990). Alternative models for
conditional stock volatility. \emph{Journal of Econometrics}, 45(1--2),
267--290.

PwC (2024). \emph{DORA: Driving Digital Operational Resilience in the EU
Financial Sector}. PwC Risk Advisory Series.

Rabiner, L. R. (1989). A tutorial on hidden Markov models and selected
applications in speech recognition. \emph{Proceedings of the IEEE},
77(2), 257--286.

Reinhart, C. M., \& Rogoff, K. S. (2009). \emph{This Time Is Different:
Eight Centuries of Financial Folly}. Princeton University Press.

Resnick, S., \& Stărică, C. (1997). Smoothing the Hill estimator.
\emph{Advances in Applied Probability}, 29(1), 271--293.

Schwert, G. W. (1989). Why does stock market volatility change over
time? \emph{Journal of Finance}, 44(5), 1115--1153.

Sims, C. A. (1980). Macroeconomics and reality. \emph{Econometrica},
48(1), 1--48.

Sims, C. A., \& Zha, T. (2006). Were there regime switches in U.S.
monetary policy? \emph{American Economic Review}, 96(1), 54--81.

Stock, J. H., \& Yogo, M. (2005). Testing for weak instruments in linear
IV regression. In Andrews, D. W. K., \& Stock, J. H. (Eds.),
\emph{Identification and Inference for Econometric Models: Essays in
Honor of Thomas Rothenberg}, 80--108. Cambridge University Press.

Taleb, N. N. (2012). \emph{Antifragile: Things That Gain from Disorder}.
Random House.

Theil, H. (1953). \emph{Repeated Least Squares Applied to Complete
Equation Systems}. The Hague: Central Planning Bureau (mimeographed).

White, H. (1980). A heteroskedasticity-consistent covariance matrix
estimator and a direct test for heteroskedasticity. \emph{Econometrica},
48(4), 817--838.

\end{document}
